\chapter{Reproducing the Results}\label{r03:app:repro}

All commands run from the workspace root unless noted. \code{.venv-rvv} is the workspace's
Python environment (at the workspace root) with lbmpy and \ps{}'s dependencies.

\section{pystencils tests on the laptop}\label{r03:sec:repro-ps}

\begin{shcode}
cd submodules/pystencils
PYTHONPATH=src ../../.venv-rvv/bin/python -m pytest -q -o addopts="" \
    -p no:cacheprovider tests/nbackend tests/codegen
# a81dc547: 435 passed, 15 skipped (13 need RVV hardware, 1 a GPU, 1 islpy)
\end{shcode}

On the board the RVV test files run natively
(\file{tests/nbackend/test_rvv_target.py}, \file{test_rvv_codegen.py},
\file{test_rvv_execution.py}, \file{test_vectorization.py}); the session scripts
\ghws{misc/x60-lab/queue_e2e_cg.sh} and \ghws{misc/x60-lab/queue_lab3.sh} show the calls.

\section{Generating lab kernels}\label{r03:sec:repro-gen}

The first argument is the output directory, the others are kernel names:
\begin{shcode}
PYTHONPATH=submodules/pystencils/src .venv-rvv/bin/python misc/x60-lab/lab/gen.py \
    misc/x60-lab/lab/src "f32-rvv-cse+aa+bpf=128+vf2+cg" "f32-rvv-cse+aa+bpf=128+vf2+cg+nomacro"
PYTHONPATH=submodules/pystencils/src .venv-rvv/bin/python misc/x60-lab/lab/gensplit.py \
    misc/x60-lab/lab/src f32-rvvsplit-by2gm5gp3 f64-rvvsplit-by1gm5gp3
\end{shcode}

\paragraph{Lab kernel names.} A generated lab kernel's name, which is also its file name in
\file{misc/x60-lab/lab/src/}, has the form \code{<precision>-<target>-<base>[+<part>...]}
(\ghws{misc/x60-lab/README.md}):
\begin{description}
  \item[\code{<precision>}] \code{f32} or \code{f64}.
  \item[\code{<target>}] \code{rvv}: RVV kernel from \ps{}; \code{generic}: plain C++ kernel.
        A \code{generic} kernel whose name contains \code{scalar} is compiled without vector
        instructions (\code{-march=rv64gc_zicbop}); all others with
        \code{-march=rv64gcv_zvl256b_zicbop_zicboz}.
  \item[\code{<base>}] \code{cse}: lbmpy's global CSE on. Any other base turns it off: the
        generator compares the base with \code{cse} exactly (\ghws{misc/x60-lab/lab/gen.py},
        function \code{emit}). This matters for \code{cse_scalar}, which is meant as ``CSE, no
        vector instructions''. The plain file \file{f32-generic-cse_scalar.cpp} is a copy of
        \file{f32-generic-cse.cpp} (its first line names that kernel), so it has CSE. But the
        files generated under a \code{cse_scalar+...} name, such as
        \file{f32-generic-cse_scalar+pf=128.cpp} and \file{f32-generic-cse_scalar+aa.cpp}, have no
        global CSE: 4 instead of 36 CSE temporaries in the pull kernel, and a checksum that
        differs from the CSE kernel's in the seventh digit. This was found while writing this
        report; \cref{r03:tab:pf} is corrected accordingly.
  \item[\code{+aa}] AA streaming: the file contains an even and an odd kernel; run with
        \code{AA=1}.
  \item[\code{+bpf=D}] software prefetch \code{D} bytes ahead, by the \ps{} option of
        \cref{r03:sec:ps-prefetch}.
  \item[\code{+pf=D}] software prefetch \code{D} bytes ahead, inserted by \code{edits.py} into the
        generated source (before the \ps{} option existed).
  \item[\code{+vf2}] a label only: generated with the new instruction selection
        (\cref{r03:sec:ps-selector}).
  \item[\code{+cg}] a label only: generated with row pointers and \code{in_register} in \ps{}
        (\cref{r03:sec:ps-rowptr}).
  \item[\code{+rp}, \code{+opq}] the same two changes, applied by \code{edits.py} to the
        generated source: row pointers, and constants hidden from the compiler.
  \item[\code{+nomacro}] the density/velocity stores and their prefetches removed: the kernel
        without output.
\end{description}
For example, \code{f32-rvv-cse+aa+bpf=128+vf2+cg} is the f32 RVV AA kernel with CSE, prefetch
128~bytes ahead, the new instruction selection, row pointers and constants in registers. The
generated split kernels are named \code{<precision>-rvvsplit-by<BY>gm<GM>gp<GP>[+pf=D][+nomacro]}
(\cref{r03:sec:wl-split}). Hand kernels are the files in \file{misc/x60-lab/lab/hand/}.

\section{QEMU check of a lab kernel}\label{r03:sec:repro-qemu}

Static cross build of the lab driver with one kernel, run under QEMU at $16^3$ for 4 timed
and 2 warm-up steps (AA kernels and split kernels need \code{AA=1}):
\begin{shcode}
L=misc/x60-lab/lab; INC=submodules/pystencils/src/pystencils/include
V="f64-rvv-cse+aa+bpf=128+vf2"          # a file pair in $L/src
riscv64-linux-gnu-g++ -O2 -march=rv64gc -DT=double -c $L/driver.cpp -o driver.o
riscv64-linux-gnu-g++ -DNDEBUG -std=c++17 -O3 -march=rv64gcv_zvl256b_zicbop \
    -I$INC -c "$L/src/$V.cpp" -o k.o
riscv64-linux-gnu-g++ -O2 -march=rv64gc -c "$L/src/$V.wrap.cpp" -o w.o
riscv64-linux-gnu-g++ -static driver.o w.o k.o -o k.bin
AA=1 qemu-riscv64 -cpu rv64,v=true,vlen=256,zicbop=true ./k.bin 16 4 2
\end{shcode}
Scalar kernels (\code{generic} and \code{scalar} in the name) are compiled with
\code{-march=rv64gc}; f32 kernels need \code{-DT=float} for the driver. The output line ends
with the three checksums of \cref{r03:sec:lab}: \code{energy=} (the sum of $|u|^2$),
\code{mass=} and \code{wsum=} (the weighted PDF sum).

\section{waLBerla on the laptop (x86)}\label{r03:sec:repro-x86}

The laptop builds were configured like the board builds of
\ghws{scripts/walberla_bench_all.sh}, with Ninja. The command below reproduces the options
recorded in the CMake cache of the build used for \cref{r03:tab:ver-app}. sweepgen needs
pystencils-sfg 0.1.1 on the \code{PYTHONPATH}; here it was installed with
\code{pip install --target}:
\begin{shcode}
export PYTHONPATH=<pystencils-sfg 0.1.1 dir>:$PWD/submodules/pystencils/src:$PWD/submodules/walberla/sweepgen/src
cmake -G Ninja -S submodules/walberla -B wb -DCMAKE_BUILD_TYPE=Release \
    -DWALBERLA_DOUBLE_ACCURACY=ON -DWALBERLA_ENABLE_V8CORE=ON -DWALBERLA_ENABLE_SWEEPGEN=ON \
    -DWALBERLA_SWEEPGEN_MANAGED_VENV=OFF -DPython_EXECUTABLE=$PWD/.venv-rvv/bin/python \
    -DWALBERLA_OPTIMIZE_FOR_LOCALHOST=ON -DWALBERLA_BUILD_WITH_MPI=ON \
    -DWALBERLA_BUILD_BENCHMARKS=ON -DWALBERLA_BUILD_TESTS=OFF
ninja -C wb SweepgenLbmCPU_explicit_cse_aa SweepgenLbmCPU_explicit_cse_aa_split
SWEEPGEN_LBM_OUTPUT=last mpirun -np 8 wb/apps/benchmarks/SweepgenLbmCPU/SweepgenLbmCPU_explicit_cse_aa_split 32 20 2
\end{shcode}
A second build directory with \code{-DWALBERLA_BUILD_TESTS=ON -DWALBERLA_BUILD_BENCHMARKS=OFF}
runs the tests: \code{ninja TestLbmScenarios_CXX TestHaloExchange TestCommunicationHiding} and
\code{ctest -R "LbmScenarios|HaloExchange|CommunicationHiding"}.

\begin{pitfall}[Generated sources are not rebuilt by themselves]
Neither Ninja nor Make tracks the Python sources of sweepgen, lbmpy or \ps{}. After changing
them, delete the generated files of the affected targets, e.g.
\file{wb/apps/benchmarks/SweepgenLbmCPU/_gen/<target>/gen/SweepgenLbmCPU.cpp} and the
matching \file{.hpp}, so that the next build regenerates them. The session scripts do this explicitly.
\end{pitfall}

\section{On the board}\label{r03:sec:repro-board}

\begin{shcode}
# one build directory, all default targets, one core, best of 3
bash scripts/walberla_lbm_bench.sh 64 200 3          # WALBERLA_BUILD=<dir>, TARGETS="..."
# a lab kernel; AA kernels need AA=1
cd misc/x60-lab/lab
RUNENV=AA=1 bash run.sh "f32-rvv-cse+aa+bpf=128+vf2"
RUNENV=AA=1 bash run.sh f32-rvvsplit-by2gm5gp3
# a session script, detached, so that it survives a dropped connection
screen -dmS e2e-final bash misc/x60-lab/queue_e2e_final.sh
\end{shcode}
How jobs run on the shared board is described in \cref{r03:sec:hw}. One more pitfall: the
board's clock runs about 1.7~days behind the laptop's, so files unpacked from an archive made
on the laptop carry modification times in the future. Make then rebuilds every dependent
target on every build until the board's clock catches up. The later session scripts reset the
times of the unpacked files with \code{touch -c}.
